%! Matrices Modeling Contexts — Unit 4, Lesson 4.14 — Homework
\documentclass[10pt]{article}
\usepackage{precalculus-article}
\usepackage{precalculus-boxes}
\newcommand{\TallMath}[1]{$\displaystyle #1\rule[-1.4em]{0pt}{3.2em}$}

\begin{document}

\pageheader{Unit 4, Lesson 4.14}{Homework}
\namedateperiod

\vspace{0.10in}

\begin{notesbox}{Practice}
\begin{enumerate}[leftmargin=1.5em, itemsep=10pt]

    \item A company tracks employees in two roles: Developer~(D) and Manager~(M).
    Each year, 90\% of Developers remain Developers and 10\% become Managers;
    80\% of Managers remain Managers and 20\% become Developers.
    Write the transition matrix $T$ (columns = ``from'' state).

    \vspace{6pt}
    $T = \begin{bmatrix}\blank{1.8cm} & \blank{1.8cm} \\ \blank{1.8cm} & \blank{1.8cm}\end{bmatrix}$

    \vspace{8pt}
    \item The initial state is $\vec{s}_0 = \begin{bmatrix}200\\50\end{bmatrix}$ (200 Developers, 50 Managers).
    Compute $\vec{s}_1$ and $\vec{s}_2$.

    \vspace{6pt}
    $\vec{s}_1 = \begin{bmatrix}\blank{5cm}\\\blank{5cm}\end{bmatrix} = \begin{bmatrix}\blank{2cm}\\\blank{2cm}\end{bmatrix}$

    \vspace{6pt}
    $\vec{s}_2 = \begin{bmatrix}\blank{5cm}\\\blank{5cm}\end{bmatrix} = \begin{bmatrix}\blank{2cm}\\\blank{2cm}\end{bmatrix}$

    \vspace{8pt}
    \item There are always 250 total employees. Set up and solve the steady-state equation to find the
    long-run number of Developers and Managers.

    \writelines{4}

    Steady state: Developers $\approx$ \blank{3cm} \quad Managers $\approx$ \blank{3cm}

    \vspace{8pt}
    \item The inverse of the matrix from problem 1 is approximately
    $T^{-1} \approx \begin{bmatrix}1.143 & 0.114 \\ -0.143 & 1.286\end{bmatrix}$.
    If $\vec{s}_1 = \begin{bmatrix}190\\60\end{bmatrix}$, use $T^{-1}$ to estimate $\vec{s}_0$.

    \vspace{6pt}
    $T^{-1}\vec{s}_1 \approx \begin{bmatrix}\blank{6cm}\\\blank{6cm}\end{bmatrix} = \begin{bmatrix}\blank{2cm}\\\blank{2cm}\end{bmatrix}$

    \vspace{8pt}
    \item \textbf{True or False} with justification: In a transition model,
    the \emph{columns} of $T$ always sum to~1.

    \vspace{4pt}
    Circle one: \quad \textbf{True} \quad \textbf{False}

    \vspace{4pt}
    Justify: \writeline

    \vspace{8pt}
    \item \textbf{AP-Style Free Response (Short).}
    A school cafeteria offers two options. Each day, 70\% of students who chose
    Option~A choose it again, and 60\% of students who chose Option~B choose it again.
    On Monday, 150 students chose A and 100 chose B.

    \begin{enumerate}[label=(\alph*), leftmargin=1.5em, itemsep=6pt]
        \item Write the transition matrix $T$.

        \vspace{4pt}
        $T = \begin{bmatrix}\blank{1.8cm} & \blank{1.8cm} \\ \blank{1.8cm} & \blank{1.8cm}\end{bmatrix}$

        \item Find Tuesday's state vector $\vec{s}_1$.

        \vspace{4pt}
        $\vec{s}_1 = \begin{bmatrix}\blank{5cm}\\\blank{5cm}\end{bmatrix} = \begin{bmatrix}\blank{2cm}\\\blank{2cm}\end{bmatrix}$

        \item What does the steady state represent in context? Answer in one or two sentences.

        \writeline
        \writeline
    \end{enumerate}

\end{enumerate}
\end{notesbox}

\end{document}
